\section{Shell Gain and Half-Derivative Branch}
\label{sec:shell-gain-half-derivative-branch}

The shell-gain branch continues the lifted-band problem after the kernel has
exposed the lower-shell derivative bill. Its positive-forward aim is to turn
active-window structure into a small coefficient. If the active shell only feels
a controlled portion of the lower-prefix strain during the relevant terminal
window, then the missing half derivative might be paid dynamically rather than
pointwise at every time.

A target active-shell coefficient has the form
\[
  A_j(t)=\frac{\text{lower-shell strain seen by }u_j}{\nu 2^{2j}}.
\]
The desired estimate is not merely that \(A_j\) is integrable. The desired
estimate is that \(A_j\) is small on the windows where the \(j\)-th shell could
contribute to terminal growth:
\[
  A_j(t)\leq \varepsilon_j,
  \qquad \sum_j \varepsilon_j<\infty.
\]
If this were available, the active shell could absorb the lower-prefix forcing
into viscosity.

The branch is a live candidate because energy gives time-integrated control. A terminal
cascade has to concentrate in time as well as scale. One hopes that the bad shell
cannot see enough lower-frequency strain for long enough to maintain a breakdown
channel. This turns the branch from a static dyadic estimate into a time-window
estimate.

The stopping point is the difference between integrated budget and active-window
control. The energy law pays
\[
  \int_0^T \|\nabla u(t)\|_{L^2}^2\,dt,
\]
while the active shell needs smallness of the coefficient it sees on its own
critical window. The missing estimate is a pointwise or window-local
coefficient bound at every active scale; the square-integrated global budget
does not supply that estimate.

This produces the half-derivative residue. The lifted-band gain reduced the
scale separation problem. The shell-gain program tries to pay the remaining
factor through time localization. The branch stalls when the time localization
requires more control than the energy budget supplies.

The residue is important for the final paper because it is a terminal survivor
pattern: precise enough to name, too strong to delete by the available estimate,
and naturally tied to a possible terminal branch. The survivor can appear as
active-window lower strain, a near-band forcing term, or a retained shell
coefficient.

In the CM-exit program, this branch presses the Field side most naturally. If the
terminal object has Pack and Part, then the question becomes whether the field
can still be read coherently at a positive scale despite the active-window
coefficient. A failure of positive-scale readout is a Field failure. If the
active-window coefficient instead shows that the selected packet has no positive
carrier, the failure moves to Pack. If it shows that the same pressure-viscosity
participation has broken, the failure moves to Part.

The shell-gain branch is not used as the positive closure. It isolates one
kind of terminal residue and sends it to the class-membership test: Field after
Pack and Part survive, or Part/Field when the selected witness loses carrier or
same-equation participation first.
